% !tp-example Basic usage
\documentclass{article}
\begin{document}
\wcTpFlipCard{
  \begin{description}
    \item[{recto}] This is the content \textbf{recto...}
    \item[{verso}] ...and here is the content \textbf{verso.}
  \end{description}
}
\end{document}

% !tp-example Hearts suit
\documentclass{article}
\begin{document}
Click any card to turn it over. You can also focus a card with Tab and press Enter or Space.


\wcTpCluster{
  \wcTpFlipCard[button-position=none,style=--tp-flip-card-padding: 0,aria-label=Ace of hearts]{
    \begin{description}
      \item[{recto}] \rawhtml|<tp-icon src="/src/components/card/cards/hearts/ha.svg" size="100%" style="display: block" role="img" aria-label="Ace of hearts"></tp-icon>|
      \item[{verso}] \rawhtml|<tp-center intrinsic="" style="height: 100%; justify-content: center"><tp-icon name="logo-tp" size="10em" role="img" aria-label="tp-components logo"></tp-icon></tp-center>|
    \end{description}
  }
  \wcTpFlipCard[button-position=none,style=--tp-flip-card-padding: 0,aria-label=Two of hearts]{
    \begin{description}
      \item[{recto}] \rawhtml|<tp-icon src="/src/components/card/cards/hearts/h2.svg" size="100%" style="display: block" role="img" aria-label="Two of hearts"></tp-icon>|
      \item[{verso}] \rawhtml|<tp-center intrinsic="" style="height: 100%; justify-content: center"><tp-icon name="logo-tp" size="10em" role="img" aria-label="tp-components logo"></tp-icon></tp-center>|
    \end{description}
  }
  \wcTpFlipCard[button-position=none,style=--tp-flip-card-padding: 0,aria-label=Three of hearts]{
    \begin{description}
      \item[{recto}] \rawhtml|<tp-icon src="/src/components/card/cards/hearts/h3.svg" size="100%" style="display: block" role="img" aria-label="Three of hearts"></tp-icon>|
      \item[{verso}] \rawhtml|<tp-center intrinsic="" style="height: 100%; justify-content: center"><tp-icon name="logo-tp" size="10em" role="img" aria-label="tp-components logo"></tp-icon></tp-center>|
    \end{description}
  }
  \wcTpFlipCard[button-position=none,style=--tp-flip-card-padding: 0,aria-label=Four of hearts]{
    \begin{description}
      \item[{recto}] \rawhtml|<tp-icon src="/src/components/card/cards/hearts/h4.svg" size="100%" style="display: block" role="img" aria-label="Four of hearts"></tp-icon>|
      \item[{verso}] \rawhtml|<tp-center intrinsic="" style="height: 100%; justify-content: center"><tp-icon name="logo-tp" size="10em" role="img" aria-label="tp-components logo"></tp-icon></tp-center>|
    \end{description}
  }
  \wcTpFlipCard[button-position=none,style=--tp-flip-card-padding: 0,aria-label=Five of hearts]{
    \begin{description}
      \item[{recto}] \rawhtml|<tp-icon src="/src/components/card/cards/hearts/h5.svg" size="100%" style="display: block" role="img" aria-label="Five of hearts"></tp-icon>|
      \item[{verso}] \rawhtml|<tp-center intrinsic="" style="height: 100%; justify-content: center"><tp-icon name="logo-tp" size="10em" role="img" aria-label="tp-components logo"></tp-icon></tp-center>|
    \end{description}
  }
  \wcTpFlipCard[button-position=none,style=--tp-flip-card-padding: 0,aria-label=Six of hearts]{
    \begin{description}
      \item[{recto}] \rawhtml|<tp-icon src="/src/components/card/cards/hearts/h6.svg" size="100%" style="display: block" role="img" aria-label="Six of hearts"></tp-icon>|
      \item[{verso}] \rawhtml|<tp-center intrinsic="" style="height: 100%; justify-content: center"><tp-icon name="logo-tp" size="10em" role="img" aria-label="tp-components logo"></tp-icon></tp-center>|
    \end{description}
  }
  \wcTpFlipCard[button-position=none,style=--tp-flip-card-padding: 0,aria-label=Seven of hearts]{
    \begin{description}
      \item[{recto}] \rawhtml|<tp-icon src="/src/components/card/cards/hearts/h7.svg" size="100%" style="display: block" role="img" aria-label="Seven of hearts"></tp-icon>|
      \item[{verso}] \rawhtml|<tp-center intrinsic="" style="height: 100%; justify-content: center"><tp-icon name="logo-tp" size="10em" role="img" aria-label="tp-components logo"></tp-icon></tp-center>|
    \end{description}
  }
  \wcTpFlipCard[button-position=none,style=--tp-flip-card-padding: 0,aria-label=Eight of hearts]{
    \begin{description}
      \item[{recto}] \rawhtml|<tp-icon src="/src/components/card/cards/hearts/h8.svg" size="100%" style="display: block" role="img" aria-label="Eight of hearts"></tp-icon>|
      \item[{verso}] \rawhtml|<tp-center intrinsic="" style="height: 100%; justify-content: center"><tp-icon name="logo-tp" size="10em" role="img" aria-label="tp-components logo"></tp-icon></tp-center>|
    \end{description}
  }
  \wcTpFlipCard[button-position=none,style=--tp-flip-card-padding: 0,aria-label=Nine of hearts]{
    \begin{description}
      \item[{recto}] \rawhtml|<tp-icon src="/src/components/card/cards/hearts/h9.svg" size="100%" style="display: block" role="img" aria-label="Nine of hearts"></tp-icon>|
      \item[{verso}] \rawhtml|<tp-center intrinsic="" style="height: 100%; justify-content: center"><tp-icon name="logo-tp" size="10em" role="img" aria-label="tp-components logo"></tp-icon></tp-center>|
    \end{description}
  }
  \wcTpFlipCard[button-position=none,style=--tp-flip-card-padding: 0,aria-label=Ten of hearts]{
    \begin{description}
      \item[{recto}] \rawhtml|<tp-icon src="/src/components/card/cards/hearts/h10.svg" size="100%" style="display: block" role="img" aria-label="Ten of hearts"></tp-icon>|
      \item[{verso}] \rawhtml|<tp-center intrinsic="" style="height: 100%; justify-content: center"><tp-icon name="logo-tp" size="10em" role="img" aria-label="tp-components logo"></tp-icon></tp-center>|
    \end{description}
  }
  \wcTpFlipCard[button-position=none,style=--tp-flip-card-padding: 0,aria-label=Jack of hearts]{
    \begin{description}
      \item[{recto}] \rawhtml|<tp-icon src="/src/components/card/cards/hearts/hj.svg" size="100%" style="display: block" role="img" aria-label="Jack of hearts"></tp-icon>|
      \item[{verso}] \rawhtml|<tp-center intrinsic="" style="height: 100%; justify-content: center"><tp-icon name="logo-tp" size="10em" role="img" aria-label="tp-components logo"></tp-icon></tp-center>|
    \end{description}
  }
  \wcTpFlipCard[button-position=none,style=--tp-flip-card-padding: 0,aria-label=Queen of hearts]{
    \begin{description}
      \item[{recto}] \rawhtml|<tp-icon src="/src/components/card/cards/hearts/hq.svg" size="100%" style="display: block" role="img" aria-label="Queen of hearts"></tp-icon>|
      \item[{verso}] \rawhtml|<tp-center intrinsic="" style="height: 100%; justify-content: center"><tp-icon name="logo-tp" size="10em" role="img" aria-label="tp-components logo"></tp-icon></tp-center>|
    \end{description}
  }
  \wcTpFlipCard[button-position=none,style=--tp-flip-card-padding: 0,aria-label=King of hearts]{
    \begin{description}
      \item[{recto}] \rawhtml|<tp-icon src="/src/components/card/cards/hearts/hk.svg" size="100%" style="display: block" role="img" aria-label="King of hearts"></tp-icon>|
      \item[{verso}] \rawhtml|<tp-center intrinsic="" style="height: 100%; justify-content: center"><tp-icon name="logo-tp" size="10em" role="img" aria-label="tp-components logo"></tp-icon></tp-center>|
    \end{description}
  }
}
\end{document}
